\documentclass[a4paper,12pt]{article}
\usepackage[margin=2.5cm]{geometry}
\usepackage{graphicx}
\usepackage{booktabs}
\usepackage{amsmath}
\usepackage{amssymb}
\usepackage{caption}
\usepackage{subcaption}
\usepackage{hyperref}
\usepackage{xcolor}
\usepackage{tabularx}

\title{\textbf{Result 1: Average End-to-End Delay vs Node Density}\\
\large ACLIF --- Autonomous Cross-Layer Intelligence Framework}
\author{B.~Tharuni~Sri~Sai\\
\small School of Computer Science and Engineering, VIT-AP University, Amaravati, India\\
\small ORCID: 0009-0004-9561-8985}
\date{\today}

\begin{document}
\maketitle
\tableofcontents
\newpage

\section{Experiment Overview}
This document presents the detailed analysis for \textbf{Result 1}: Average End-to-End Delay vs Node Density,
as part of the ACLIF cross-layer IoT optimization framework evaluation.

\section{Result Figure}
\begin{figure}[htbp]
  \centering
  \includegraphics[width=0.92\textwidth]{../figures/fig_delay_nodes.png}
  \caption{Average End-to-End Delay vs Node Density. Results obtained from NS3 v3.38 simulations over a
  100-node heterogeneous IoT network ($200\times200$~m$^2$), averaged over 10 runs
  with 95\% confidence interval shading.}
  \label{fig:delay_vs_nodes}
\end{figure}

\section{Analysis and Discussion}
This experiment evaluates the average end-to-end delay as the number of sensor nodes
increases from 20 to 100 in a $200\times200$~m$^2$ deployment area. ACLIF consistently
achieves the lowest delay across all node densities. At 100 nodes, ACLIF achieves
91.5~ms, outperforming the best baseline IPSO at 118.2~ms. This 22.6\% improvement
stems from the DQN agent's ability to simultaneously adapt MAC contention window,
transmit power, and routing metric, reducing queuing at congested relay nodes.
The joint optimization identifies that at high density, load-balanced routing
($a_\mathrm{route}=2$) steers flows away from bottleneck nodes, while contention
window expansion directly suppresses collision-induced retransmissions.

\subsubsection*{Key Metric}
\begin{equation}
D_\mathrm{e2e}(\mathcal{P}) = \sum_{j=1}^{m-1}\left(
  \frac{d_{n_j,n_{j+1}}}{c} + \frac{q_{n_j}(t)}{\mu_{n_j}} + \frac{k}{R_{n_j}}
\right)
\end{equation}

\subsubsection*{Summary Statistics}
\begin{tabular}{lcc}
\toprule
Method & Delay @ 100 nodes (ms) & \% vs ACLIF \\
\midrule
E-GLBR & 148.3 & +62.1\% \\
FDRL   & 121.7 & +33.0\% \\
FACSO  & 163.9 & +79.1\% \\
EEHCT  & 135.4 & +47.9\% \\
IPSO   & 118.2 & +29.2\% \\
\textbf{ACLIF} & \textbf{91.5} & --- \\
\bottomrule
\end{tabular}

\section{Experimental Configuration}
\begin{tabular}{ll}
\toprule
\textbf{Parameter} & \textbf{Value} \\
\midrule
Simulator      & NS3 v3.38 \\
Nodes          & 100 sensor nodes \\
Area           & $200 \times 200$~m$^2$ \\
PHY standard   & IEEE 802.15.4 \\
MAC protocol   & CSMA/CA \\
Packet payload & 512 bytes \\
Packet rate    & 5 pkt/s (default) \\
Runs per config & 10 (different seeds) \\
CI             & 95\% Student-$t$ \\
\bottomrule
\end{tabular}

\section{Baselines}
ACLIF is compared against: E-GLBR \cite{Benelhouri2023}, FDRL \cite{Suresh2024},
FACSO \cite{Yang2024204}, EEHCT \cite{Chaurasiya202325941}, and IPSO \cite{Zhang2025}.

\begin{thebibliography}{9}
\bibitem{Benelhouri2023} Benelhouri et al., ``E-GLBR Genetic Algorithm Routing,'' 2023.
\bibitem{Suresh2024} Suresh et al., ``FDRL Federated Deep RL Routing,'' 2024.
\bibitem{Yang2024204} Yang et al., ``FACSO Fuzzy Adaptive Cross-Layer Scheduling,'' 2024.
\bibitem{Chaurasiya202325941} Chaurasiya et al., ``EEHCT Energy-Efficient Clustering,'' 2023.
\bibitem{Zhang2025} Zhang et al., ``IPSO-Based Multi-Hop Routing,'' 2025.
\end{thebibliography}

\end{document}
